\chapter{Conception du système}

\section{Introduction}

Ce chapitre présente les choix de conception retenus pour transformer les besoins du chapitre~2 en une architecture, des modèles de données et des scénarios de fonctionnement cohérents. Les diagrammes décrivent le \textbf{système implémenté} et distinguent les responsabilités des clients, du backend, des services internes et des acteurs humains.

\section{Architecture générale}

Le système est une architecture composée de deux clients (web et mobile) et d'un backend commun, communiquant par HTTP(S) et WebSocket, avec une base documentaire MongoDB. Le client web repose sur Next.js / React et le client mobile sur Expo / React Native. Le backend expose une API REST/JSON, coordonne l'authentification, les conversations, le moteur RAG/Gemini et les échanges temps réel avec Socket.io.

\begin{figure}[H]
\centering
\begin{tikzpicture}[
  font=\small,
  box/.style={draw=sinapsbleu, rounded corners=3pt, align=center, inner sep=7pt, fill=white, drop shadow},
  title/.style={font=\bfseries\small\color{sinapsbleu}}
]
\node[box, fill=sinapsbleu!8, minimum width=4.2cm] (fe) {\textbf{Frontend}\\ Next.js 16 / React 19\\ TypeScript, Tailwind\\ \texttt{:3000}};
\node[box, fill=sinapsbleu!8, minimum width=4.6cm, right=1.3cm of fe] (be) {\textbf{Backend}\\ Node.js / Express 5\\ Socket.io, JWT, Multer\\ \texttt{:5000}};
\node[box, fill=sinapsbleu!8, minimum width=4.2cm, below=1.35cm of fe, xshift=-2.5cm] (mobile) {\textbf{Client mobile}\\ Expo / React Native\\ HTTP + Socket.io};
\node[box, fill=grisclair, minimum width=3.6cm, below=1.5cm of be] (db) {\textbf{MongoDB}\\ User, Agent,\\ Conversation, Message};
\node[box, fill=grisclair, minimum width=3.6cm, left=1.2cm of db] (ai) {\textbf{IA}\\ Retriever TF-IDF\\ + API Gemini};
\node[box, fill=grisclair, minimum width=3.4cm, right=1.2cm of db] (ext) {\textbf{Externes}\\ Google OAuth\\ fichiers \texttt{/uploads}};

\draw[-{Latex}] (fe.east) -- node[above, font=\scriptsize, yshift=2pt]{REST / JSON} (be.west);
\draw[-{Latex}] (be.west) -- node[below, font=\scriptsize, yshift=-2pt]{Socket.io} (fe.east);
\draw[-{Latex}] (mobile.east) -- node[below, font=\scriptsize]{REST / Socket.io} (be.south west);
\draw[-{Latex}] (be.south) -- (db.north);
\draw[-{Latex}] (be.south west) -- (ai.north east);
\draw[-{Latex}] (be.south east) -- (ext.north west);
\end{tikzpicture}
\caption{Architecture logique d'Op Support IA (vue déployable en local)}
\label{fig:archi}
\end{figure}

Les flux principaux sont :

\begin{enumerate}[leftmargin=1.6em]
  \item le navigateur charge l'interface Next.js ;
  \item les appels métier passent par l'API Express ;
  \item après un message client, si la conversation est en mode IA, le backend appelle le retriever puis Gemini (ou le repli) et enregistre un message \texttt{ia} ;
  \item Socket.io notifie la salle de conversation et, pour les listes, un événement global ;
  \item les fichiers sont stockés sur disque sous \texttt{uploads/} et servis en statique ;
  \item le client peut demander une prise en charge humaine ; l'agent traite alors la conversation depuis sa file.
\end{enumerate}

\section{Architecture logicielle}

L'organisation logique repose sur quatre ensembles complémentaires :

\begin{itemize}[leftmargin=1.4em]
  \item les \textbf{clients}, web et mobile, qui présentent les interfaces de conversation, d'agent et d'administration et communiquent avec le backend ;
  \item le \textbf{backend}, qui expose l'API REST, applique les contrôles d'accès, persiste les données et coordonne les changements d'état ;
  \item les \textbf{services internes}, qui regroupent l'authentification, le retriever RAG, la génération Gemini, le téléversement et les notifications Socket.io ;
  \item la \textbf{persistance et les services externes}, constitués de MongoDB pour les entités métier, de Google OAuth pour l'identification client et de Gemini pour la génération lorsqu'une clé valide est disponible.
\end{itemize}

Cette séparation permet aux deux clients de partager les mêmes conversations, messages et règles d'autorisation sans dupliquer la logique métier. Les composants techniques restent distincts des acteurs humains identifiés au chapitre~2.

\section{Diagramme de cas d'utilisation}

\begin{figure}[H]
\centering
\resizebox{\textwidth}{!}{%
\begin{tikzpicture}[
  font=\footnotesize,
  actor/.style={align=center},
  uc/.style={draw, ellipse, minimum width=2.9cm, minimum height=0.85cm, align=center, inner sep=1pt},
  sys/.style={draw=sinapsbleu, thick, dashed, inner sep=12pt, rounded corners}
]
\node[actor] (c) at (0,0) {\textbf{Client}};
\node[actor] (a) at (0,-5.2) {\textbf{Agent}};
\node[actor] (ad) at (11.2,-2.4) {\textbf{Admin}};

\node[uc] (u1) at (4.6,1.6) {S'identifier};
\node[uc] (u2) at (4.6,0.5) {Discuter (IA)};
\node[uc] (u3) at (4.6,-0.6) {Joindre un fichier};
\node[uc] (u4) at (4.6,-1.7) {Escalader / revenir IA};
\node[uc] (u5) at (4.6,-2.8) {Clôturer + noter};

\node[uc] (u6) at (4.6,-4.2) {Se connecter};
\node[uc] (u7) at (4.6,-5.3) {Répondre (humain)};
\node[uc] (u8) at (4.6,-6.4) {Prendre une conversation};

\node[uc] (u9) at (8.6,-0.4) {Valider / rejeter agent};
\node[uc] (u10) at (8.6,-1.6) {Voir les statistiques};
\node[uc] (u11) at (8.6,-2.8) {Filtrer l'historique};
\node[uc] (u12) at (8.6,-4.2) {S'inscrire (pending)};

\node[sys, fit=(u1)(u2)(u3)(u4)(u5)(u6)(u7)(u8)(u9)(u10)(u11)(u12)] (sys) {};
\node[above] at (sys.north) {\textbf{Op Support IA}};

\draw (c) -- (u1);
\draw (c) -- (u2);
\draw (c) -- (u3);
\draw (c) -- (u4);
\draw (c) -- (u5);
\draw (a) -- (u6);
\draw (a) -- (u7);
\draw (a) -- (u8);
\draw (a) -- (u3);
\draw (ad) -- (u9);
\draw (ad) -- (u10);
\draw (ad) -- (u11);
\draw (a) -- (u12);
\end{tikzpicture}%
}
\caption{Cas d'utilisation principaux}
\label{fig:cu}
\end{figure}

\subsection{CU01 --- S'identifier (client)}

Le client fournit un \textbf{jeton Google}, vérifié avec \texttt{GOOGLE\_CLIENT\_ID}, ou un couple nom / e-mail. Le serveur trouve ou crée le profil client et délivre un JWT de rôle \texttt{client}. Une conversation non résolue est réutilisée, sinon une nouvelle conversation est créée avec \texttt{handledBy = ia} et \texttt{status = en\_cours}.

\subsection{CU02 --- Envoyer un message}

Avec une session client et une conversation ouverte, le message est persisté. Si \texttt{handledBy = ia}, le service IA produit une réponse enregistrée comme message \texttt{ia}. Les deux messages sont émis sur Socket.io.

\subsection{CU03 --- Escalader}

Le client, ou un profil autorisé sur la conversation, passe \texttt{handledBy} à \texttt{humain} et \texttt{status} à \texttt{en\_attente}. Les agents voient alors le fil dans la file.

\subsection{CU04 --- Traiter une demande (agent)}

L'agent authentifié et approuvé consulte la liste, peut être assigné, répond avec \texttt{sender = humain}, puis clôture la conversation.

\subsection{CU05 --- Administrer}

L'administrateur consulte \texttt{GET /api/stats} et \texttt{GET /api/conversations}, approuve ou supprime un agent en attente et filtre les conversations.

\section{Modèle de données}

MongoDB stocke quatre collections Mongoose. Les relations sont des références \texttt{ObjectId}.

\begin{table}[H]
\centering
\caption{Collection \texttt{users} (clients)}
\label{tab:user}
\begin{tabularx}{\textwidth}{@{}llX@{}}
\toprule
Champ & Type & Commentaire \\
\midrule
googleId & String, unique, requis & Identifiant Google ou, en flux direct, souvent l'e-mail \\
name, email & String, requis & Profil affiché \\
avatar & String & URL d'avatar \\
timestamps & Date & \texttt{createdAt}, \texttt{updatedAt} \\
\bottomrule
\end{tabularx}
\end{table}

\begin{table}[H]
\centering
\caption{Collection \texttt{agents} (opérateurs et administrateurs)}
\label{tab:agent}
\begin{tabularx}{\textwidth}{@{}llX@{}}
\toprule
Champ & Type & Commentaire \\
\midrule
name, email & String (e-mail unique) & Identité de connexion \\
password & String & Hash bcrypt \\
role & \texttt{agent} $\mid$ \texttt{admin} & Défaut : agent \\
skills & [String] & Saisies à l'inscription ; \textbf{non utilisées} pour un routage automatique \\
status & \texttt{pending} $\mid$ \texttt{approved} & Un agent non approuvé ne peut pas se connecter \\
avatar & String & Optionnel \\
\bottomrule
\end{tabularx}
\end{table}

\begin{table}[H]
\centering
\caption{Collection \texttt{conversations}}
\label{tab:conv}
\begin{tabularx}{\textwidth}{@{}llX@{}}
\toprule
Champ & Type & Commentaire \\
\midrule
client & ObjectId $\rightarrow$ User & Propriétaire \\
assignedAgent & ObjectId $\rightarrow$ Agent & Peut être nul \\
status & \texttt{en\_cours} $\mid$ \texttt{en\_attente} $\mid$ \texttt{resolu} & Cycle de vie \\
handledBy & \texttt{ia} $\mid$ \texttt{humain} & Qui répond \\
satisfaction.rating & Number 1--5 & Optionnel à la clôture \\
satisfaction.comment & String & Optionnel \\
resolvedBy, resolvedAt & String, Date & Origine et instant de résolution \\
resolutionType & Enum & Résolution confirmée par le client, par l'agent ou abandon \\
lastActivityAt & Date & Sélection de la conversation active la plus récente \\
aiAttemptCount, escalationCount & Number & Suivi des tentatives IA et des escalades \\
isTestData & Boolean & Marqueur exclu des indicateurs opérationnels \\
\bottomrule
\end{tabularx}
\end{table}

\begin{table}[H]
\centering
\caption{Collection \texttt{messages}}
\label{tab:msg}
\begin{tabularx}{\textwidth}{@{}llX@{}}
\toprule
Champ & Type & Commentaire \\
\midrule
conversation & ObjectId & Fil parent \\
sender & \texttt{client} $\mid$ \texttt{ia} $\mid$ \texttt{humain} & Provenance \\
authorName & String & Affichage \\
content & String & Texte (peut être vide si pièce jointe) \\
attachments[] & url, type, name & \texttt{image} $\mid$ \texttt{video} $\mid$ \texttt{document} $\mid$ \texttt{link} \\
quickReplies[] & id, label, action, metadata & Actions proposées après une réponse IA \\
timestamps & Date & \texttt{createdAt}, \texttt{updatedAt}, utilisés pour l'affichage du fil \\
\bottomrule
\end{tabularx}
\end{table}

\section{Diagramme de classes}

\begin{figure}[H]
\centering
\begin{tikzpicture}[
  font=\scriptsize,
  cls/.style={draw=sinapsbleu, rectangle split, rectangle split parts=2, rounded corners=2pt, align=left, fill=white, inner sep=4pt}
]
\node[cls] (user) {
  \textbf{User}\\
  googleId\\ name\\ email\\ avatar
  \nodepart{two}
  timestamps
};
\node[cls, right=0.55cm of user] (agent) {
  \textbf{Agent}\\
  name\\ email\\ password\\ role\\ skills[]\\ status\\ avatar
  \nodepart{two}
  login / signup
};
\node[cls, below=1.3cm of user] (conv) {
  \textbf{Conversation}\\
  status\\ handledBy\\ satisfaction
  \nodepart{two}
  escalate()\\ deescalate()\\ assign()\\ close()
};
\node[cls, right=1.25cm of conv] (msg) {
  \textbf{Message}\\
  sender\\ authorName\\ content\\ attachments[]
  \nodepart{two}
  ---
};

\draw[-{Latex}] (conv) -- node[pos=0.15, left, font=\tiny]{0..*} node[pos=0.85, right, font=\tiny]{1} (user);
\draw[-{Latex}] (conv.east) -- node[pos=0.15, right, font=\tiny]{0..1} node[pos=0.85, left, font=\tiny]{0..*} (agent.south);
\draw[-{Latex}] (msg) -- node[pos=0.15, above, font=\tiny]{1..*} node[pos=0.85, below, font=\tiny]{1} (conv);
\end{tikzpicture}
\caption{Diagramme de classes aligné sur les schémas Mongoose}
\label{fig:classes}
\end{figure}

\section{Diagrammes de séquence}

\subsection{Envoi d'un message client et réponse IA}

Le client transmet le message au backend, qui le persiste, interroge le retriever et le service Gemini lorsque la conversation est prise en charge par l'IA. La réponse éventuelle est enregistrée puis diffusée aux clients par Socket.io.

\begin{figure}[H]
\centering
\resizebox{\textwidth}{!}{%
\begin{tikzpicture}[font=\scriptsize, x=1.7cm, y=-0.55cm]
\foreach \x/\n in {0/Client, 2/Next.js, 4/Express, 6/MongoDB, 8/RAG+Gemini} {
  \node[font=\bfseries\scriptsize] at (\x,0) {\n};
  \draw[thin] (\x,-0.4) -- (\x,-11.2);
}
\draw[-{Latex}] (0,-1) -- node[above]{\texttt{POST /messages}} (4,-1);
\draw[-{Latex}] (4,-2) -- node[above]{insert Message client} (6,-2);
\draw[-{Latex}] (6,-3) -- (4,-3);
\draw[-{Latex}] (4,-4) -- node[above]{retrieve + generate} (8,-4);
\draw[-{Latex}] (8,-5.2) -- (4,-5.2);
\draw[-{Latex}] (4,-6.4) -- node[above]{insert Message ia} (6,-6.4);
\draw[-{Latex}] (4,-7.6) -- node[above]{Socket \texttt{message\_received}} (2,-7.6);
\draw[-{Latex}] (2,-8.8) -- (0,-8.8);
\node[align=left] at (4,-10.2) {Uniquement si \texttt{handledBy = ia}.\\Sinon, pas d'appel IA.};
\end{tikzpicture}%
}
\caption{Séquence : message client et génération IA}
\label{fig:seq-ia}
\end{figure}

\subsection{Escalade vers un humain}

Le client appelle \texttt{PATCH /conversations/:id/escalate}. Le document passe à \texttt{handledBy = humain} et \texttt{status = en\_attente}. Socket.io émet \texttt{conversation\_updated} vers la salle ; un événement global permet à l'espace agent de rafraîchir la file.

\subsection{Clôture et satisfaction}

\texttt{PATCH /conversations/:id/close} accepte une note entière entre 1 et 5. Une valeur hors intervalle, par exemple 0, n'est pas enregistrée comme satisfaction. Le statut devient \texttt{resolu} et la cause de résolution est tracée dans \texttt{resolutionType}. Les actions rapides permettent notamment de confirmer une résolution IA, de poser une nouvelle question ou de demander un agent.

\section{Cycle de vie d'une conversation}

Une conversation est créée ou réutilisée lors de l'identification du client. Elle commence avec \texttt{handledBy = ia} et \texttt{status = en\_cours}. Les messages sont alors traités par l'IA lorsque le client écrit dans une conversation ouverte. La figure~\ref{fig:conversation-state} synthétise les états et transitions décrits ci-dessous.

Lorsque le client demande une prise en charge humaine, ou lorsqu'un acteur autorisé effectue cette action, la conversation passe à \texttt{handledBy = humain} et \texttt{status = en\_attente}. Un agent approuvé peut la prendre en charge ; elle est alors éventuellement associée à \texttt{assignedAgent} et repasse à \texttt{status = en\_cours}. L'agent peut répondre, puis clôturer la conversation.

La clôture fait passer le statut à \texttt{resolu}. Une note entière de 1 à 5 et un commentaire peuvent être enregistrés ; la résolution conserve également son origine et son instant. Le retour vers l'IA est possible avant la clôture par désescalade. Les statuts et modes associés sont récapitulés dans l'annexe~\ref{tab:annexe-statuts}.

\begin{figure}[H]
\centering
\begin{tikzpicture}[
  font=\scriptsize,
  state/.style={draw=sinapsbleu, thick, rounded corners=3pt, align=center, minimum width=3.7cm, minimum height=0.85cm, fill=sinapsbleu!8},
  terminal/.style={draw=sinapsbleu, thick, circle, minimum size=0.45cm, fill=sinapsbleu!25},
  transition/.style={-{Latex}, draw=sinapsbleu, thick},
  note/.style={font=\scriptsize, align=center}
]
\node[terminal] (start) {};
\node[state, below=0.55cm of start] (ia) {EN\_COURS\\prise en charge par l'IA};
\node[state, below=1.05cm of ia] (waiting) {EN\_ATTENTE\\demande d'aide humaine};
\node[state, below=1.05cm of waiting] (human) {EN\_COURS\\prise en charge humaine};
\node[state, below=1.05cm of human] (resolved) {RESOLU};
\node[terminal, below=0.55cm of resolved] (end) {};
\node[note, right=0.55cm of ia] (iaNote) {traitement IA};
\node[note, right=0.55cm of waiting] (waitNote) {assignation à un agent};
\node[note, right=0.55cm of human] (humanNote) {retour vers l'IA};
\node[note, right=0.55cm of resolved] (ratingNote) {évaluation facultative};

\draw[transition] (start) -- (ia);
\draw[transition] (ia) -- node[left, note] {demande d'aide humaine} (waiting);
\draw[transition] (waiting) -- node[left, note] {agent assigné} (human);
\draw[transition] (human) -- node[left, note] {clôture} (resolved);
\draw[transition] (resolved) -- (end);
\draw[transition] (human.east) .. controls +(2.0,0) and +(2.0,0) .. node[right, note] {désescalade} (ia.east);
\draw[transition] (ia.west) .. controls +(-2.0,0) and +(-2.0,0) .. node[left, note] {confirmation / clôture} (resolved.west);
\end{tikzpicture}
\caption{Cycle de vie d'une conversation dans SINAPS}
\label{fig:conversation-state}
\end{figure}

\section{Authentification et gestion des droits}

Deux familles de jetons JWT signés avec \texttt{JWT\_SECRET} coexistent.

\begin{table}[H]
\centering
\caption{Politique d'authentification}
\label{tab:auth}
\small
\begin{tabularx}{\textwidth}{@{}lX@{}}
\toprule
Profil & Mécanisme \\
\midrule
Client & JWT \texttt{role=client}, 30~jours, après Google (ID token vérifié) ou identification e-mail. Stockage navigateur : \texttt{sinaps\_client}. \\
Agent / admin & JWT \texttt{role=agent$\mid$admin}, 7~jours, après e-mail + mot de passe. Refus si agent \texttt{pending}. Stockage : \texttt{sinaps\_token}. \\
Middleware & \texttt{requireAuth}, \texttt{requireAdmin}, \texttt{requireAgentOrAdmin}, \texttt{requireClientAuth}, \texttt{requireConversationAccess}, \texttt{requireSenderAuth}, \texttt{requireAnySession}. \\
\bottomrule
\end{tabularx}
\end{table}

Les règles de conception sont les suivantes :

\begin{itemize}[leftmargin=1.4em]
  \item un client ne lit et ne modifie que \textbf{sa} conversation ;
  \item \texttt{sender = humain} exige un JWT agent/admin ;
  \item \texttt{sender = ia} est refusé s'il vient d'une requête externe ;
  \item l'upload exige une session client ou agent.
\end{itemize}

En production, le serveur refuse de démarrer sans \texttt{JWT\_SECRET}. La mise en œuvre détaillée des routes et des middleware est présentée au chapitre~4.

\section{Conception du RAG}

\subsection{Retriever}

La base \texttt{data/knowledgeBase.js} contient des paires \texttt{\{ question, answer \}}. Chaque document est normalisé et tokenisé, puis indexé en TF et pondéré par IDF. La requête et les documents sont comparés par \textbf{similarité cosinus}. Les extraits dont le score est strictement supérieur à $0{,}35$ sont retenus, avec un maximum de deux documents (\texttt{topK = 2}).

Ce n'est \textbf{pas} un RAG vectoriel : aucun embedding ni base vectorielle n'est utilisé. Il s'agit d'un RAG \og lexical \fg{} adapté à une petite FAQ de démonstration.

\subsection{Génération et repli}

Si une clé Gemini valide est présente, les extraits pertinents sont injectés dans un prompt français afin de produire une réponse courte et fondée sur la base de connaissances. Les modèles Gemini configurés sont essayés en cascade. Si la clé est absente, si l'appel échoue ou si aucun document pertinent n'est disponible, le système renvoie la meilleure réponse locale lorsqu'elle existe, sinon un message invitant à demander l'aide d'un agent humain.

\subsection{Escalade}

L'escalade n'est \textbf{pas automatique} selon un score de confiance. Elle est \textbf{décidée par l'utilisateur} ou par un acteur autorisé sur la conversation via l'action prévue à cet effet. Le mécanisme RAG peut conduire à proposer cette aide humaine lorsqu'aucune réponse pertinente n'est disponible, mais le changement de prise en charge reste une action explicite.

\section{Conclusion}

Ce chapitre a présenté l'architecture générale et logicielle, les cas d'utilisation, le modèle de données, les séquences, le cycle de vie des conversations, l'authentification et la conception du RAG. Le chapitre suivant décrit la réalisation concrète de ces choix dans les applications web, mobile et backend.
